\paragraph{\textbf{Proposition 3.12.}}
\SGARefTarget{sga3:xx:3.12}
\emph{Let \((G,T,\alpha)\) be an elementary \(S\)-system and let
\(w\in N^\times(S)\).  Put}
\[
  \Omega_0=\Omega\cap\intaut(w^{-1})(\Omega).
\]
\emph{Let \(d\) be the function on \(\Omega\) defined by}
\[
  d\bigl(\exp(Y)\cdot t\cdot\exp(X)\bigr)
    =\alpha(t)^{-1}+XY.
\]
\emph{Then \(\Omega_0=\Omega_d\), and, for
\(\exp(Y)\cdot t\cdot\exp(X)\in\Omega_0(S')\), one has the following
formula, where
\(z=d(\exp(Y)\cdot t\cdot\exp(X))\):}
\begin{equation}
\tag{\(\star\)}\label{sga3:xx:proposition:3-12:equation:star}
  \intaut(w)\bigl(\exp(Y)\cdot t\cdot\exp(X)\bigr)
  =
  \exp\bigl(z^{-1}a_\alpha(w)^{-1}X\bigr)\cdot
  t\alpha^\ast(z)\cdot
  \exp\bigl(z^{-1}a_\alpha(w)Y\bigr).
\end{equation}
\emph{Moreover,}
\[
  d\circ\intaut(w)=d^{-1}.
\]

Indeed, one immediately has\footnote{The editors corrected the
original by interchanging \(a_\alpha(w)\) and
\(a_\alpha(w)^{-1}\), and supplied the details of the proof of
\(d\circ\intaut(w)=d^{-1}\).}
\[
\begin{aligned}
  \intaut(w)\bigl(\exp(Y)\cdot t\cdot\exp(X)\bigr)
   &=
   \exp\bigl(a_\alpha(w)Y\bigr)\cdot
   t\alpha^\ast\bigl(\alpha(t)^{-1}\bigr)\cdot
   \exp\bigl(a_\alpha(w)^{-1}X\bigr)\\
   &=
   \exp\bigl(a_\alpha(w)Y\bigr)\cdot
   \exp\bigl(\alpha(t)a_\alpha(w)^{-1}X\bigr)\cdot
   t\alpha^\ast\bigl(\alpha(t)^{-1}\bigr).
\end{aligned}
\]
By \SGARefLink{sga3:xx:2.1}{2.1}, this is a section of \(\Omega\) if and only if
\(1+\alpha(t)XY\) is invertible, which proves
\(\Omega_0=\Omega_d\).  Applying formula \textup{(F)} of the same
paragraph then gives \((\star)\) by an immediate computation.
Finally, \((\star)\) gives
\[
\begin{aligned}
  (d\circ\intaut(w))
    \bigl(\exp(Y)\cdot t\cdot\exp(X)\bigr)
   &=
   \alpha\bigl(t\alpha^\ast(z)\bigr)^{-1}
     +z^{-2}XY\\
   &=z^{-2}\bigl(\alpha(t)^{-1}+XY\bigr)
    =z^{-1},
\end{aligned}
\]
which proves the last assertion.

\smallskip
\noindent\textbf{N.B.}\ The function \(d\) is independent of the
choice of \(w\).

\section{The Isomorphism Theorem}

\paragraph{\textbf{Theorem 4.1.}}
\SGARefTarget{sga3:xx:4.1}
\emph{Let \(S\) be a scheme and let \(q\in\mathbb Z\), \(q>0\), be
such that \(x\mapsto x^q\) is an endomorphism of
\(\mathbb G_{a,S}\).  Let \((G,T,\alpha)\) and
\((G',T',\alpha')\) be two elementary \(S\)-systems.  Let}
\[
  h\colon(\mathfrak g^\alpha)^{\otimes q}
      \xrightarrow{\ \sim\ }\mathfrak g'^{\,\alpha'},
  \qquad
  h^\vee\colon(\mathfrak g^{-\alpha})^{\otimes q}
      \xrightarrow{\ \sim\ }\mathfrak g'^{-\alpha'}
\]
\emph{be mutually contragredient isomorphisms.  Let
\(f_T\colon T\to T'\) be a morphism of \(S\)-groups making the
following diagram commute:}

\begin{center}
  \begin{tikzcd}[column sep=huge,row sep=large]
    \mathbb G_{m,S}
      \arrow[r,"q"]
      \arrow[d,"\alpha^\ast"']
      & \mathbb G_{m,S}
        \arrow[d,"\alpha'^\ast"] \\
    T
      \arrow[r,"f_T"]
      \arrow[d,"\alpha"']
      & T'
        \arrow[d,"\alpha'"] \\
    \mathbb G_{m,S}
      \arrow[r,"q"]
      & \mathbb G_{m,S}
  \end{tikzcd}
\end{center}

\emph{There exists a unique morphism of \(S\)-groups
\(f\colon G\to G'\) extending \(f_T\) and satisfying}
\[
  f\bigl(\exp(X)\bigr)=\exp\bigl(h(X^q)\bigr)
\]
\emph{for every \(S'\to S\) and
\(X\in W(\mathfrak g^\alpha)(S')\).  Moreover, this morphism also
satisfies}
\[
  f\bigl(\exp(Y)\bigr)=\exp\bigl(h^\vee(Y^q)\bigr),
  \qquad
  f\bigl(w_\alpha(Z)\bigr)
    =w_{\alpha'}\bigl(h(Z^q)\bigr)
\]
\emph{for every \(S'\to S\),
\(Y\in\Gamma(S',\mathfrak g^{-\alpha})\), and
\(Z\in\Gamma(S',\mathfrak g^\alpha)^\times\).}

If \(f\colon G\to G'\) extends \(f_T\), then
\[
  f\circ\alpha^\ast=(\alpha'^\ast)^q.
\]
If, in addition, \(f\) satisfies the second displayed condition in
the statement, then it satisfies the other two conditions as well by
\SGARefLink{sga3:xx:3.10}{3.10}.  It follows that \(f\) is determined on \(\Omega\) by
\[
  f\bigl(\exp(Y)t\exp(X)\bigr)
  =
  \exp\bigl(h^\vee(Y^q)\bigr)\,
  f_T(t)\,
  \exp\bigl(h(X^q)\bigr).
\]
Since \(\Omega\) is schematically dense in \(G\), this already proves
the uniqueness of \(f\).  To prove existence, it is enough, by
\SGARefLink{sga3:xviii:prop:2.3}{Exposé XVIII, 2.3}, to show that the preceding formula defines a
``generically multiplicative'' morphism from \(\Omega\) to \(G'\).
By \SGARefLink{sga3:xx:2.4}{2.4}, this amounts to checking
\[
  \alpha'\circ f=\alpha^q,
\]
which follows from the fact that \(f\) extends \(f_T\).

\paragraph{\textbf{Scholium 4.2.}}
\SGARefTarget{sga3:xx:4.2}
\SGARefLink{sga3:xx:4.1}{Theorem 4.1} may also be interpreted as follows.  Let
\(\mathcal E\) be the category of elementary \(S\)-systems, and let
\(\mathcal D\) be the category of tuples
\[
  \bigl(
    \mathbb G_{m,S}\xrightarrow{\alpha^\ast}T
    \xrightarrow{\alpha}\mathbb G_{m,S},\,L
  \bigr),
\]
where \(T\) is a torus, \(\alpha\) and \(\alpha^\ast\) are group
morphisms such that \(\alpha\circ\alpha^\ast=2\), and \(L\) is an
invertible \(\mathcal O_S\)-module.  The reader may specify the
morphisms in the two categories under consideration.  Define a
functor
\[
\begin{aligned}
  \mathcal E&\longrightarrow\mathcal D,\\
  (G,T,\alpha)&\longmapsto
  \bigl(
    \mathbb G_{m,S}\xrightarrow{\alpha^\ast}T
    \xrightarrow{\alpha}\mathbb G_{m,S},\,
    \mathfrak g^\alpha
  \bigr).
\end{aligned}
\]
The preceding theorem says that this functor is fully faithful.  In
fact, it is an equivalence of categories, as will be seen in the next
section.  We already have:

\paragraph{\textbf{Corollary 4.3.}}
\SGARefTarget{sga3:xx:4.3}
\emph{If \(q=1\) and \(f_T\) is an isomorphism, then \(f\) is an
isomorphism.}

\paragraph{\textbf{Corollary 4.4.}}
\SGARefTarget{sga3:xx:4.4}
\emph{If \(q=1\) and \(f_T\) is faithfully flat with kernel \(Q\)
\textup{(cf.\ \SGARefLink{sga3:ix:proposition:2:7}{Exposé IX, 2.7})}, then \(f\) is faithfully flat and
quasi-compact with kernel \(Q\), and hence identifies \(G'\) with
\(G/Q\).}

Indeed, if \(f_T\) is faithfully flat with kernel \(Q\), then
\[
  Q=\Ker(f_T)
  \subset\Ker(\alpha'\circ f_T)
  =\Ker(\alpha).
\]
Introduce the elementary \(S\)-system
\((G/Q,T/Q,\alpha/Q)\) of \SGARefLink{sga3:xx:2.13}{2.13}.  By \SGARefLink{sga3:xx:2.14}{2.14}, we are reduced to proving
that \(f/Q\) induces an isomorphism from \(G/Q\) onto \(G'\), which
follows immediately from \SGARefLink{sga3:xx:4.3}{4.3}.

